%
% 公立はこだて未来大学卒業研究中間報告書[全コース対応版]
%
%         ファイル名："sample.tex"
%
% !TEX encoding = UTF-8
% !TEX program = lualatex
\documentclass[11pt]{ltjarticle}
\usepackage{funinfosys}
\usepackage{url}
\usepackage{graphicx}
\author{%
1022179 西野理希\\指導教員 : イアン・フランク
}
%\course{Information Systems Course}
%\course{Advanced ICT Course} %% 高度ICTコースの場合はこちらを使用
%\course{Information Design Course} %% 情報デザインコースの場合はこちらを使用
\course{Complex Systems Course} %% 複雑系科学コースの場合はこちらを使用
%\course{Intelligent Systems Course} %% 知能システムコースの場合はこちらを使用
\title{作るか、選ぶか?\\教員による提示方法が教育動画の価値認識に与える影響}
\etitle{Create or Curate?\\ How Instructor Framing Shapes the Value of Educational Video}
\eauthor{NISHINO Riki}
\abstract{大学教育では教育動画の活用が広がっている一方，教員が独自に動画を制作するには時間や技術的な負担が伴う．本研究では，既存の動画を選定するキュレーションに着目し，教員による動画の提示方法が学習者の価値認識に与える影響を検討する．大学生を対象としたオンライン実験を計画し，参加者を「教員が制作した動画」「教員が厳選した動画」「由来を説明しない動画」の3条件に割り当てる．参加者は各条件で共通の教育動画3本を視聴し，動画ごとに7段階のアンケートと理解度テストに回答する．主要評価指標は，教材としての適切性，授業への採用価値，他の学生への推奨評価から構成する教育的価値とし，提示条件間の差を分析する．現在は生成AIの文章生成を題材とした動画の脚本初稿を作成しており，今後は動画と実験手順の準備を進める．}
\keywords{教育動画，キュレーション，提示方法，価値認識}
\eabstract{Educational videos are widely used in higher education, but producing them can require substantial time and technical skills. This study examines whether instructor framing affects students' perceived value of educational videos, focusing on the distinction between creating and curating materials. An online experiment with university students is planned. Participants will be randomly assigned to one of three conditions: videos described as created by the instructor, videos described as carefully curated by the instructor, or videos presented without information about their origin. All participants will watch the same three videos and complete seven-point rating questionnaires and comprehension tests after each video. The primary outcome is perceived educational value, assessed through appropriateness as course material, value for classroom use, and willingness to recommend the video to other students. A draft script on generative AI has been prepared.}
\ekeywords{Educational Videos, Curation, Framing, Perceived Value}
\begin{document}
\maketitle
%\vspace*{-.5cm}

\section{Background and Objectives}

In recent years, educational videos have become more widely used in higher education alongside the expansion of online learning environments. Noetel et al.\ \cite{noetel2021} reported that educational videos contribute to improved learning outcomes in higher education. However, producing educational videos requires time and technical skills and can place a burden on instructors who create their own materials. Curating existing videos by selecting and using them for educational purposes is therefore another way to prepare teaching materials.

In the use of educational videos, creating materials and selecting existing materials represent different forms of instructor involvement. However, it remains insufficiently clear how information about such involvement affects students' evaluations of teaching materials.

This study aims to examine how instructor framing of educational videos affects students' perceived value. Specifically, three conditions will be compared in which identical videos are described as ``created by the instructor,'' ``selected by the instructor,'' or presented without information about their origin. Students' ratings of perceived educational value, usefulness, credibility, and related aspects will be compared. By keeping the video content constant, the study seeks to identify the influence of differences in descriptions of instructor involvement on students' evaluations.

\section{Relevance to the Complex Systems Course}

The Complex Systems Course aims to develop the ability to analyze regularities in natural and social phenomena and apply this understanding to real-world problems. Through a comparative experiment, this study examines the relationship between the framing of instructor involvement in educational videos and students' perceived value. By quantitatively analyzing ratings obtained through manipulation of the framing conditions, it seeks to investigate how information about instructor involvement affects student evaluations and to provide insights into the presentation of teaching materials in educational settings.

\section{Related Work}

Regarding the effectiveness of educational videos, Noetel et al.\ \cite{noetel2021} reported that their use in higher education contributes to improved learning outcomes. Liesa-Or\'us et al.\ \cite{liesa2023} discussed challenges related to university instructors' digital competence. Since producing educational videos also requires such competence, it is meaningful to consider ways of using existing teaching materials.

As an approach to using existing materials, Deschaine and Sharma \cite{deschaine2015} characterized digital curation as selecting, organizing, and sharing learning resources in accordance with educational purposes. Flintoff et al.\ \cite{flintoff2014} also discussed the importance of digital curation in education and research. These studies suggest that instructors' roles include selecting existing materials as well as producing them.

Meanwhile, Hovland and Weiss \cite{hovland1951} showed that source credibility affects the reception of information. Ramlatchan and Watson \cite{ramlatchan2020} reported that the presentation of online lecture videos affects students' perceptions of instructor credibility and psychological immediacy. These findings suggest that evaluations may be related not only to information content but also to how it is presented and how its source is perceived.

Building on these findings, this study compares students' evaluations of identical educational videos described either as ``created by the instructor'' or ``selected by the instructor.'' It thereby examines the effects of descriptions of instructor involvement on perceived value, rather than the effects of how the teaching materials were actually created or selected.

\section{Research Methods}

\subsection{Experimental Design}

An online experiment with university students is planned. Participants will be randomly assigned to one of three conditions---Create, Curate, or Neutral---to compare different ways of framing instructor involvement in videos. The wording presented in each condition will be as follows:

\begin{itemize}
\item \textbf{Create:} ``This video was created by the course instructor for use in this course.''
\item \textbf{Curate:} ``This video was selected by the course instructor for use in this course.''
\item \textbf{Neutral:} ``This video will be used in this course.''
\end{itemize}

Each participant will watch three educational videos on AI literacy, each approximately three minutes long, under their assigned condition. The same videos will be used in all conditions, and presentation order will be balanced as far as practicable. Only the description of instructor involvement will vary between conditions; the videos themselves will remain unchanged.

\subsection{Evaluation Measures}

After each video, participants will complete a seven-point rating questionnaire and a content-comprehension test. The primary outcome is the educational value of the video as course material as subjectively perceived by students. Three items will assess suitability as course material, value for adoption in a course, and willingness to recommend the video to other students. The mean of these three ratings is being considered as the primary measure.

Secondary measures will assess usefulness for learning, credibility of the explanation, satisfaction, and motivation to learn, using one item for each. In addition, each video will be followed by two or three multiple-choice comprehension questions to provide supplementary information about understanding of the content.

After all videos have been evaluated, participants will be asked what they remember being told about the videos' origin and how strongly they believed that explanation. These questions will help establish how participants understood the framing information.

\subsection{Analysis Methods}

In the primary analysis, each participant's perceived-educational-value scores will be averaged across the three videos to obtain a participant-level score. These scores will be used to compare the Create and Curate conditions, while comparisons with the Neutral condition will be considered secondary.

Depending on differences in ratings across videos, the number of participants, and other factors, a linear mixed-effects model may also be considered for additional analysis \cite{baayen2008}. Results will be interpreted with attention to differences in ratings between conditions and the uncertainty surrounding those differences.

\subsection{Research Limitations and Ethical Considerations}

The number of participants and the range of videos used may be limited, which would constrain the generalizability of the findings. The validity of the study-specific evaluation items and possible differences in quality and difficulty among videos must also be considered.

Furthermore, depending on the condition, participants may receive an explanation of a video's origin that differs from its actual production history. Ethical considerations will therefore need to be addressed, including information provided to participants, informed consent, and explanation after the experiment.

\section{Current Progress and Future Plans}

\subsection{Development of Stimulus Videos}

Preparations are underway for the educational videos to be used in the experiment. Generative AI is being considered as the subject area, and an initial script has been prepared for a video titled ``How Does Generative AI Produce Responses?''

The video is intended for university students who are beginners in AI and explains the basic mechanism by which generative AI generates text. It is planned as an approximately three-minute video combining diagrams, animation, and narration. The visual content and presentation have been designed as five scenes comprising 19 shots. The main topics are tokens as units of text processing and the repeated prediction and addition of the next token during text generation.

\subsection{Future Plans}

The next step is to complete one educational video and consider a trial run incorporating the questionnaire and experimental procedure. This trial run is intended to check the entire sequence, from presentation of the framing condition and video viewing to comprehension questions and questionnaire responses, and to identify issues to address before the main experiment.

Subsequently, the remaining educational videos will be prepared, and participant recruitment and experimental procedures will be specified in greater detail. In the main experiment, the collected ratings will be compared across framing conditions to analyze how differences in the explanation of instructor involvement affect students' perceived value.

\section{Conclusion}

This study examines how instructor framing of educational videos affects students' perceived value. Current work includes designing a comparative experiment with Create, Curate, and Neutral conditions, considering evaluation measures and analysis methods, and preparing an educational video script. The next steps are to check the experimental procedure and questionnaire through a trial run using one completed video, prepare the remaining videos, and work toward conducting the online experiment. The study aims to clarify how differences in descriptions of instructor involvement affect students' evaluations of educational videos.

\begin{thebibliography}{99}\small
\bibitem{noetel2021}
Noetel, M., et al. (2021). Video improves learning in higher education: A systematic review. \textit{Review of Educational Research}, 91(2), 204--236. \url{https://doi.org/10.3102/0034654321990713}
\bibitem{liesa2023}
Liesa-Or\'us, M., Lozano Blasco, R., \& Arce-Romeral, L. (2023). Digital competence in university lecturers: A meta-analysis of teaching challenges. \textit{Education Sciences}, 13(5), 508. \url{https://doi.org/10.3390/educsci13050508}
\bibitem{deschaine2015}
Deschaine, M. E., \& Sharma, S. A. (2015). The five Cs of digital curation: Supporting twenty-first-century teaching and learning. \textit{InSight: A Journal of Scholarly Teaching}, 10, 19--24. \url{https://doi.org/10.46504/10201501de}
\bibitem{flintoff2014}
Flintoff, K., Mellow, P., \& Clark, K. P. (2014). Digital curation: Opportunities for learning, teaching, research and professional development. In \textit{Proceedings of the 23rd Annual Teaching Learning Forum}. University of Western Australia.
\bibitem{hovland1951}
Hovland, C. I., \& Weiss, W. (1951). The influence of source credibility on communication effectiveness. \textit{Public Opinion Quarterly}, 15(4), 635--650. \url{https://doi.org/10.1086/266350}
\bibitem{ramlatchan2020}
Ramlatchan, M., \& Watson, G. S. (2020). Enhancing instructor credibility and immediacy in online multimedia designs. \textit{Educational Technology Research and Development}, 68(1), 511--528. \url{https://doi.org/10.1007/s11423-019-09714-y}
\bibitem{baayen2008}
Baayen, R. H., Davidson, D. J., \& Bates, D. M. (2008). Mixed-effects modeling with crossed random effects for subjects and items. \textit{Journal of Memory and Language}, 59(4), 390--412. \url{https://doi.org/10.1016/j.jml.2007.12.005}
\end{thebibliography}
\end{document}
%
%
% EOF
